\section{Discussion and Limitations}\label{sec:discussion}

Checking each label against a model that never saw it is a simple way to clean a small training set for a linear classifier. In the tested settings, one fixed filter runs without being told the noise rate, or whether the labels are noisy at all, and under noise it leaves a much cleaner training set. It does not restore clean-label accuracy: its mean test accuracy stays below the clean-label reference in all four noisy settings. It also discards correctly labeled samples, and on breast cancer at the higher rate, many mislabeled ones survive it.

\paragraph{Limitations and next steps.}
We studied two small datasets, synthetic symmetric noise at two rates, one filtering pass with fixed $K$ and $C$, and one untuned MLP, which we did not combine with the filter. We tested the filter only where the clean-label reference is very accurate. Removing samples at a fixed $C$ also changes the relative strength of regularization, and we did not try stronger regularization without filtering. Without random-drop, oracle or in-sample filter controls, we cannot separate the effects of surviving mislabeled samples, lost correctly labeled samples and a smaller training set. We ran neither co-teaching~\citep{han2018coteaching} nor confident learning~\citep{northcutt2021confident} and claim no advantage over either. Next steps are these missing runs, class-dependent noise, other rates, and softer or repeated filtering.
